\documentclass[10pt]{article}
\usepackage[T1]{fontenc}
\usepackage[utf8]{inputenc}
\usepackage[french]{babel}
\usepackage{amsmath,amssymb,mathrsfs,mathpazo}
\usepackage{geometry}
\geometry{a4paper,margin=2cm}
\usepackage{graphicx,array,tabularx,tabularray,colortbl,multirow,multicol,booktabs,enumitem,pifont,ulem,fancybox,diagbox,cancel}
\usepackage{pstricks,pstricks-add,pst-plot,pst-tree,pst-node,pst-text,pst-func,pst-math,pst-bspline,pst-3dplot,pst-eucl,pst-solides3d}
\let\pscalculate\undefined
\usepackage{pst-bezier}
\usepackage{tikz,pgfplots,tkz-tab}
\usetikzlibrary{arrows,arrows.meta,patterns,calc,positioning,trees,decorations.markings,shapes.geometric,shapes.misc,shapes.arrows,intersections,decorations.pathmorphing,decorations.pathreplacing,backgrounds,shadows,babel}
\pgfplotsset{compat=1.18}
\usepgfplotslibrary{fillbetween}
\usepackage[np]{numprint}
\usepackage{listings,makecell,pas-tableur,twoopt}
\let\barre\undefined
\usepackage[pstricks]{tablvar}
\definecolor{darkgreen}{RGB}{0,100,0}
\providecommand{\red}{\color{red}}
\providecommand{\blue}{\color{blue}}
\providecommand{\green}{\color{green}}
\providecommand{\black}{\color{black}}
\providecommand{\white}{\color{white}}
\providecommand{\gray}{\color{gray}}
\providecommand{\cyan}{\color{cyan}}
\providecommand{\magenta}{\color{magenta}}
\providecommand{\yellow}{\color{yellow}}
\providecommand{\violet}{\color{violet}}
\providecommand{\orange}{\color{orange}}
\providecommand{\euro}{\texteuro}
\providecommand{\EUR}{\texteuro}
\providecommand{\ieme}{\textsuperscript{e}}
\providecommand{\degres}{\ensuremath{^\circ}}
\providecommand{\degree}{\ensuremath{^\circ}}
\providecommand{\vv}[1]{\overrightarrow{#1}}
\providecommand{\Coord}[1]{#1}
\providecommand{\decosix}{}
\providecommand{\starredbullet}{\ensuremath{\bullet}}
\pagestyle{empty}
\setlength{\parindent}{0pt}
\usepackage{adjustbox}\DeclareUnicodeCharacter{2236}{:}\DeclareUnicodeCharacter{2190}{\ensuremath{\leftarrow}}
\providecommand{\R}{\mathbb{R}}
\providecommand{\N}{\mathbb{N}}
\providecommand{\Z}{\mathbb{Z}}
\providecommand{\Q}{\mathbb{Q}}
\providecommand{\D}{\mathbb{D}}
\providecommand{\C}{\mathbb{C}}
\providecommand{\e}{{\mathrm{e}}}
\providecommand{\vect}[1]{\overrightarrow{#1}}
\providecommand{\vectt}[1]{\overrightarrow{\mathrm{#1}}}
\providecommand{\vvt}[1]{\overrightarrow{\mathrm{#1}}}
\providecommand{\np}[2][]{{#2~#1}}
\providecommand{\vv}[1]{{\overrightarrow{#1}}}
\renewcommand{\d}{{\mathrm{d}}}
\providecommand{\degres}{{^\circ}}
\providecommand{\degre}{{^\circ}}
\providecommand{\blue}{\color{blue}}
\providecommand{\red}{\color{red}}
\providecommand{\enskip}{\quad}
\providecommand{\llbracket}{\lbrack\!\lbrack}
\providecommand{\rrbracket}{\rbrack\!\rbrack}
\providecommand{\euro}{{\text{€}}}
\providecommand{\no}{{\text{n°}}}
\providecommand{\interval}[2]{{[#1;#2]}}
\providecommand{\intervFO}[2]{{[#1;#2[}}
\providecommand{\dotfill}{\ldots}

\providecommand{\ds}{\displaystyle}
\providecommand{\textding}[1]{#1}
\providecommand{\Oij}{\(\mathrm{O};\vec{\imath},\vec{\jmath}\)}
\providecommand{\Oijk}{\(\mathrm{O};\vec{\imath},\vec{\jmath},\vec{k}\)}
\providecommand{\Ouv}{\(\mathrm{O};\vec{u},\vec{v}\)}
\providecommand{\intervFF}[2]{\ensuremath{[#1;#2]}}
\renewcommand{\euro}{\ifmmode\text{\texteuro}\else\texteuro\fi}
\DeclareUnicodeCharacter{00B0}{\ifmmode^{\circ}\else\textdegree\fi}
\DeclareUnicodeCharacter{20AC}{\ifmmode\text{\texteuro}\else\texteuro\fi}
% Énoncés : Éducation nationale. Transcription et corrigés : APMEP, Denis Vergès et contributeurs. Découpage et classement : CoopMaths, Sébastien Lozano, Éric Elter et contributeurs. Adaptation de la présentation : Maths Explorer.
% CC BY-SA : https://coopmaths.fr/www/about/
% https://coopmaths.fr/alea/static/bac/2025/tex/bac_2025_09_sujet1_polynesie_1_cor.tex
\begin{document}
\noindent\textit{APMEP / CoopMaths. CC BY-SA. Présentation adaptée par Maths Explorer.}\par\medskip
\begin{list}{\textbullet}{En France il y a deux formules pour obtenir le permis de conduire:}
			\item Suivre à partir de 15 ans une formation de conduite accompagnée pendant 2 ans ;
		\item Suivre la formation classique (sans conduite accompagnée) à partir de 17 ans.
	\end{list}

	

	En France actuellement, parmi les jeunes qui suivent une formation au permis de con-duire, 16~\% choisissent la formation de conduite accompagnée, et parmi eux, 74,7~\% réussissent l'examen de conduite dès leur première tentative.
	En suivant la formation classique, le taux de réussite dès la première tentative est seulement de 56,8~\%.
	On choisit au hasard un jeune français qui a déjà passé l'examen de conduite et on considère les évènements $A$ et $R$ suivants :

	\begin{list}{\textbullet}{}
		\item $A$ : \og le jeune a suivi la formation de conduite accompagnée\fg{};
		\item $R$ : \og le jeune a eu le permis dès sa première tentative \fg{}.
	\end{list}







	\subsection*{Partie A}

		\begin{enumerate}
			\item On dresse un arbre de probabilités modélisant cette situation.
			
\begin{center}
\bigskip
\psset{nodesepB=4pt,levelsep=35mm,treesep=10mm}
\pstree[treemode=R,nrot=:U]{\TR{}}
{\pstree[nodesepA=4pt]{\TR{$A$}\naput{\small $0,16$}}
		{\TR{$R$}\naput{\small $0,747$}
		\TR{$\overline{R}$}\nbput{\small $1-0,747=0,253$}
		}
\pstree[nodesepA=4pt]{\TR{$\overline{A}$}\nbput{\small $1-0,16=0,84$}}
		{\TR{$R$}\naput{\small $0,568$}
		\TR{$\overline{R}$}\nbput{\small $1-0,568 = 0,432$}
		}
	}
	\bigskip
\end{center}			
			
			\item \begin{enumerate}
				\item 
D'après la loi des probabilités totales:

$\begin{aligned}
P(R)
& = P(A \cap R) + P\left(\overline{ A} \cap R\right) 
= P(A) \times P_A(R) + P\left(\overline{A}\right) \times P_{\overline{A}}(R)\\
& = 0,16 \times 0,747 + 0,84 \times 0,568 = \np{0,11952} + \np{0,47712}\\
& = \np{0,59664}
\end{aligned}$
				
			\end{enumerate}

Dans la suite, on gardera la valeur $0,597$ arrondie à $10^{-3}$ près.

			\begin{enumerate}[resume]
				\item 
Une probabilité de $0,597$ correspond à un pourcentage de $59,7\,\%$.

Le pourcentage de réussite lors du premier passage du permis est de $59,7\,\%$, quel que soit le mode d'apprentissage.				
			\end{enumerate}

			\item On choisit un jeune ayant eu son permis dès sa première tentative. 

La probabilité qu'il ait suivi la formation de conduite accompagnée est:

$P_{R}(A)= \dfrac{P\left (A\cap R\right )}{P\left (R\right )} \approx \dfrac{0,16\times 0,747}{0,597} \approx 0,200$

			\item On cherche la proportion $x$ de jeunes suivant la formation de conduite accompagnée pour que le taux de réussite global (quelle que soit la formation choisie) dès la première tentative à l'examen de conduite dépasse $70~\%$.
			
			L'arbre précédent devient:
			
\begin{center}
\bigskip
\psset{nodesepB=4pt,levelsep=35mm,treesep=10mm}
\pstree[treemode=R]{\TR{}}
{\pstree[nodesepA=4pt,nrot=:U]{\TR{$A$}\naput{\small $x$}}
		{\TR{$R$}\naput{\small $0,747$}
		\TR{$\overline{R}$}\nbput{\small $0,253$}
		}
\pstree[nodesepA=4pt,nrot=:U]{\TR{$\overline{A}$}\nbput{\small $1-x$}}
		{\TR{$R$}\naput{\small $0,568$}
		\TR{$\overline{R}$}\nbput{\small $0,432$}
		}
	}
	\bigskip
\end{center}		

Le taux de réussite devient alors:

$P(R) =x\times  0,747 + (1-x)\times 0,568 = 0,747x + 0,568 - 0,568x = 0,179x + 0,568$

$P(R) > 0,70
\iff 0,179x + 0,568 > 0,7 
\iff 0,179x > 0,132 
\iff x > \dfrac{0,132}{0,179}$

Or $\dfrac{0,132}{0,179} \approx \np{0,737}$.

Il faudrait donc qu'au moins $73,7\,\%$ des jeunes choisissent la conduite accompagnée pour dépasser un taux de réussite de 70\,\%.
			
		\end{enumerate}

	\subsection*{Partie B}
		Une auto-école présente pour la première fois à l'examen de conduite 10 candidats qui ont suivi la formation de conduite accompagnée. On modélise le fait de passer les examens de conduite par des épreuves aléatoires indépendantes.



		On note $X$ la variable aléatoire donnant le nombre de ces 10 candidats qui auront leur permis dès la première tentative.


		\begin{enumerate}
			\item 
L'expérience élémentaire consiste à choisir un candidat; il y a deux possibilités: il est reçu, avec une probabilité $p=0,747$, ou il ne l'est pas.

On réalise 10 fois cette expérience et on suppose que les épreuves aléatoires sont indépendantes.

Donc  la variable aléatoire $X$ donnant le nombre de ces 10 candidats qui auront leur permis dès la première tentative suit la loi binomiale de paramètres $n=10$ et $p=0,747$.

			\item 
$P(X \geqslant 6) = 1-  P(X \leqslant 5)$

À la calculatrice on trouve  $P(X \leqslant 5) \approx 0,082$ donc $P(X \geqslant 6) \approx 0,918$.

On peut donc estimer à $91,8\,\%$ le pourcentage de chances que sur 10 candidats, au moins 6 obtiennent le permis dès leur première tentative.

			\item 
$E(X)=np=10\times 0,747=7,47$ et $V(X)=np\left (1-p\right )=10\times 0,747 \times 0,253 \approx 1,890$

			\item Il y a aussi 40 candidats qui n'ont pas suivi la formation de conduite accompagnée et qui se présentent pour la première fois à l'examen de conduite. De la même manière, on note $Y$ la variable aléatoire qui donne le nombre de ces candidats qui auront le permis à la première tentative. On admet que $Y$ est indépendante de la variable $X$ et qu'en fait $E(Y)=22,53$ et $V(Y)=9,81$.



			On note alors $Z$ la variable aléatoire comptant le nombre total de candidats (parmi les 50) qui auront le permis de conduire dès la première tentative.



			\begin{enumerate}
				\item 
D'après le texte: $Z=X+Y$ 

D'après la linéarité de l'espérance: $E(Z)=E(X)+E(Y) = 7,47+ 22,53 = 30$.

Les variables $X$ et $Y$ sont indépendantes, donc on peut utiliser l'additivité de la variance: $V(Z)=V(X)+V(Y) \approx 1,89 + 9,81 = 11,7$.

	\item 
	
La variable aléatoire $Z$ a pour espérance $E(Z)$ et pour variance $V(Z)$ donc, d'après l'inégalité de Bienaymé-Tchebychev, on a:

pour tout $\delta\in\, \left ]0\;;\;+\infty\strut \right [$, $P\left ( \left |Z-E(Z)\strut \right | \geqslant \delta \right ) \leqslant \dfrac{V(Z)}{\delta^2}$.

$E(Z)=30$ et $V(Z)=11,7$ donc  l'inégalité de Bienaymé-Tchebychev devient:

pour tout $\delta\in\, \left ]0\;;\;+\infty\strut \right [$, $P\left ( \left |Z-30\strut \right | \geqslant \delta \right ) \leqslant \dfrac{11,7}{\delta^2}$.

On cherche la probabilité qu'il y ait moins de 20  ou plus de 40 candidats reçus.

$ Z\leqslant 20 \iff Z-30 \leqslant -10$ et $Z \geqslant 40 \iff Z-30 \geqslant 10$

$\left ( (Z-30 \leqslant -10) \text{ ou } (Z-30 \geqslant 10) \strut \right )\iff \left |Z-30 \strut \right | \geqslant 10$

Pour $\delta=10$, l'inégalité de Bienaymé-Tchebychev devient:

$P\left ( \left |Z-30\strut \right | \geqslant 10 \right ) \leqslant \dfrac{11,7}{10^2}$
donc
$P\left ( \left |Z-30\strut \right | \geqslant 10 \right ) \leqslant 0,117$

donc la probabilité qu'il y ait au plus 20 ou au moins 40 candidats qui aient leur permis dès la première tentative est inférieure à 0,12.
			\end{enumerate}
		\end{enumerate}

\bigskip

\end{document}
